\begin{summarybox}
\textbf{\textcolor{CSummary}{一句话核心}}
\textbf{正实数函数 $x\ln x$ 是凸函数，其 Jensen 不等式给出了多个点函数值均值与均值点函数值的大小关系。}

\medskip
\textbf{\textcolor{CSummary}{使用条件}}
\begin{itemize}[leftmargin=*, itemsep=0.5em, topsep=0.4em]
    \item \textbf{条件 1（变量正实数）：} 所有变量 $x_i > 0$，确保 $\ln x_i$ 有意义且凸性成立。
    \item \textbf{条件 2（凸函数可用）：} 确认 $f(x)=x\ln x$ 的形式，且结论基于其凸性。
\end{itemize}

\medskip
\textbf{\textcolor{CSummary}{关键提醒}}
\begin{itemize}[leftmargin=*, itemsep=0.5em, topsep=0.4em]
    \item \textbf{易错点（定义域检查）：} 使用前务必检查每个 $x_i$ 是否为正数，否则结论失效。
    \item \textbf{检查项（取等条件）：} 完成放缩后，务必写出等号成立的条件（所有变量相等），这是使用严格凸函数不等式的完整步骤。
\end{itemize}
\end{summarybox}
